\begin{definitionbox}{Definition (Associativity)}
\begin{definition}[Associativity]
\label{def:associativity}
\[
\forall x,y,z,\ x\mathbin{O}(y\mathbin{O}z)=(x\mathbin{O}y)\mathbin{O}z.
\]
\end{definition}
\end{definitionbox}

\begin{theorembox}{Theorem (General Associativity)}
\begin{theorem}[General Associativity]
\label{thm:general-associativity}
If $O$ is associative, then every bracketing of
$x_1\mathbin{O}\cdots\mathbin{O}x_n$ has the same value.
\end{theorem}
\end{theorembox}

\begin{dependencies}
\begin{itemize}
  \item \hyperref[def:associativity]{Associativity}
\end{itemize}
\end{dependencies}

\begin{definitionbox}{Definition (Commutativity)}
\begin{definition}[Commutativity]
\label{def:commutativity}
\[
\forall x,y,\ x\mathbin{O}y=y\mathbin{O}x.
\]
\end{definition}
\end{definitionbox}

\begin{theorembox}{Theorem (General Commutativity)}
\begin{theorem}[General Commutativity]
\label{thm:general-commutativity}
If $O$ is associative and commutative, then any finite reordering of
$x_1\mathbin{O}\cdots\mathbin{O}x_n$ has the same value.
\end{theorem}
\end{theorembox}

\begin{dependencies}
\begin{itemize}
  \item \hyperref[def:associativity]{Associativity}
  \item \hyperref[def:commutativity]{Commutativity}
\end{itemize}
\end{dependencies}

\begin{definitionbox}{Definition (Left Distributivity)}
\begin{definition}[Left Distributivity]
\label{def:left-distributivity}
The operation $\otimes$ is left distributive over $\oplus$ if
\[
x\otimes(y\oplus z)=(x\otimes y)\oplus(x\otimes z),
\]
\end{definition}
\end{definitionbox}

\begin{definitionbox}{Definition (Right Distributivity)}
\begin{definition}[Right Distributivity]
\label{def:right-distributivity}
The operation $\otimes$ is right distributive over $\oplus$ if
\[
(y\oplus z)\otimes x=(y\otimes x)\oplus(z\otimes x),
\]
\end{definition}
\end{definitionbox}

\begin{definitionbox}{Definition (Two-Sided Distributivity)}
\begin{definition}[Two-Sided Distributivity]
\label{def:two-sided-distributivity}
The operation $\otimes$ is two-sided distributive over $\oplus$ if it is both
left distributive and right distributive.
\end{definition}
\end{definitionbox}

\begin{definitionbox}{Definition (Left Cancellation)}
\begin{definition}[Left Cancellation]
\label{def:left-cancellation}
An operation has left cancellation if
\[
x\mathbin{O}y=x\mathbin{O}z \Longrightarrow y=z,
\]
\end{definition}
\end{definitionbox}

\begin{definitionbox}{Definition (Right Cancellation)}
\begin{definition}[Right Cancellation]
\label{def:right-cancellation}
An operation has right cancellation if
\[
y\mathbin{O}x=z\mathbin{O}x \Longrightarrow y=z.
\]
\end{definition}
\end{definitionbox}

\begin{definitionbox}{Definition (Two-Sided Cancellation)}
\begin{definition}[Two-Sided Cancellation]
\label{def:two-sided-cancellation}
An operation has two-sided cancellation if it has both left cancellation and
right cancellation.
\end{definition}
\end{definitionbox}

\begin{propositionbox}{Proposition (Invertibility Implies Left Cancellation)}
\begin{proposition}[Invertibility Implies Left Cancellation]
\label{prop:invertibility-implies-left-cancellation}
In an associative operation with identity, invertible elements satisfy left
cancellation.
\end{proposition}
\end{propositionbox}

\begin{dependencies}
\begin{itemize}
  \item \hyperref[def:associativity]{Associativity}
  \item \hyperref[def:two-sided-identity]{Two-Sided Identity}
  \item \hyperref[def:two-sided-inverse]{Two-Sided Inverse}
\end{itemize}
\end{dependencies}

\begin{propositionbox}{Proposition (Invertibility Implies Cancellation)}
\begin{proposition}[Invertibility Implies Cancellation]
\label{prop:invertibility-implies-cancellation}
In an associative operation with identity, invertible elements satisfy
two-sided cancellation.
\end{proposition}
\end{propositionbox}

\begin{dependencies}
\begin{itemize}
  \item \hyperref[prop:invertibility-implies-left-cancellation]{Invertibility Implies Left Cancellation}
  \item \hyperref[prop:invertibility-implies-right-cancellation]{Invertibility Implies Right Cancellation}
  \item \hyperref[def:two-sided-cancellation]{Two-Sided Cancellation}
\end{itemize}
\end{dependencies}

\begin{propositionbox}{Proposition (Invertibility Implies Right Cancellation)}
\begin{proposition}[Invertibility Implies Right Cancellation]
\label{prop:invertibility-implies-right-cancellation}
In an associative operation with identity, invertible elements satisfy right
cancellation.
\end{proposition}
\end{propositionbox}

\begin{dependencies}
\begin{itemize}
  \item \hyperref[def:associativity]{Associativity}
  \item \hyperref[def:two-sided-identity]{Two-Sided Identity}
  \item \hyperref[def:two-sided-inverse]{Two-Sided Inverse}
\end{itemize}
\end{dependencies}
